\documentclass[9pt,a4paper]{article}

\usepackage[T1]{fontenc}
\usepackage[utf8]{inputenc}
\usepackage[cyr]{aeguill}
\usepackage[francais]{babel}

% %\renewcommand{\baselinestretch}{1.5}  % line spacing : 1,5
% \setlength{\oddsidemargin}{2,5cm}	% Marge gauche sur pages impaires
% \setlength{\evensidemargin}{2,5cm}	% Marge gauche sur pages paires
% \setlength{\textwidth}{16cm}	% Largeur de la zone de texte (17cm)
% \setlength{\topmargin}{2,5cm}	% Pas de marge en haut

% Les packages
%\usepackage[french]{babel}
\usepackage{fancybox}
\usepackage{graphics}
\usepackage{color}
\usepackage{latexsym}
\usepackage{amsmath,amsthm,amssymb}
\usepackage[plain]{fullpage} 
\usepackage{verbatim}
\usepackage{times,epsfig}
\usepackage{listings}
\usepackage{multirow}
\usepackage{amsmath,amsthm,amssymb,amscd,txfonts,bbm} % RAJOUT ICI !!!
\usepackage{graphics,epsfig} % RAJOUT ICI !!!
%\usepackage{algorithmic}
\usepackage[final]{pdfpages} 
\usepackage{fancyvrb}
\usepackage{enumerate}
\usepackage{listingsutf8}
\usepackage{tikz}
\usetikzlibrary{arrows,shapes,snakes,automata,backgrounds,petri,positioning}
\usepackage{tikz-qtree}
\usepackage{comment}


\newenvironment{qv}
{\quote\Verbatim}
{\endquote\endVerbatim}

\usepackage{lastpage}

\newcommand{\terme}[1]{\texttt{#1}}
\newcommand{\termSet}[1]{\{\texttt{#1}\}}
\newcommand{\guill}[1]{<<\,#1\,>>}
% Macro pour l'inclusion de figure  
\newcommand{\fig}[2]{%
  \begin{figure*}[hbt]
   \begin{center}
      \includegraphics[width=6cm]{#1}
  \caption{\label{#1}  #2}
  \end{center}
 \end{figure*}
}
\newcommand{\figTex}[2]{%
  \begin{figure*}[hbt]
 \centerline{\input{#1.pstex_t}}
  \caption{\label{#1}  #2}
 \end{figure*}
}

% Macro commentaires
\definecolor{turquoise}{rgb}{0.20 ,0.60, 0.60}
\definecolor{vert}{rgb}{0.40,  0.60 ,0.00}
\definecolor{violet}{rgb}{0.80 , 0.00 , 0.80}
\definecolor{bleu}{rgb}{0.20,  0.20,  0.80}
\definecolor{rose}{rgb}{1.00,  0.20,  0.40}
\def\remVG#1{\textcolor{vert}{\texttt{Virginie: #1}}}
\def\comVG#1{\begin{footnotesize}\textcolor{vert}{\texttt{Precision de Virginie: #1}}\end{footnotesize}}

\lstset{basicstyle=\footnotesize,showstringspaces=false,language=XML,breaklines=false,columns=fullflexible,frame=shadowbox, captionpos=b}

\newenvironment{solution}[0]{
~\\ \color{red} \textbf{Solution :}\color{black}}{%
}
%\excludecomment{solution}


\parindent=0pt           
                             
\usepackage{url}

\newtheorem{Exercice}{Definition}

\lstset{language=XML,basicstyle=\footnotesize,
numberstyle=\footnotesize, 
breaklines=true
showspaces=false,         
showstringspaces=false,
showtabs=false,
frame=single,
tabsize=4}   


\begin{document}

\title{Examen NFE204 - Session 1}
\date{04 février 2014}

\maketitle

\begin{center}
{\bf \Large Consignes}\\
\begin{tabular}{| p{15cm} |}
\hline
\\
Dur\'ee de l'examen : 2h00\\
Documents non autoris\'es;  Calculatrice autoris\'ee\\

Les exercices sont ind\'ependants les uns des autres. Merci de soigner votre \'ecriture.\\
Ce sujet contient \pageref{LastPage} pages.\\


\textit{Important}\,: nous n'évaluons pas les réponses par la syntaxe. ne vous
inquiétez pas d'utiliser un point-virgule à la place dun point d'exclamation ou
autres détails mineurs
\\
\hline
\end{tabular}
\end{center}

\section*{Première partie : XML (7 pts)}

On va s'intéresser aux recettes de cuisine et à leur représentation documentaire. Voici tout
d'abord un exemple de recette représenté en XML.

\vspace{-0.05cm} 
\lstinputlisting[inputencoding=utf8/latin1,basicstyle=\small,language=XML,breaklines=true, caption={Fichier crepes.xml}]{crepes.xml}


\subsubsection*{XPath: compréhension}

Donnez le résultat des expressions XPath suivantes, appliqué au document \texttt{crepes.xml}.
\begin{enumerate}
  \item \verb7/recipe/ingredient_lines/ingredient_line/ingredient7
  \item \verb7/recipe/ingredient_lines/ingredient_line/quantity/@unit7
  \item \verb7/recipe/catalogued_in/catalog[@name="difficulté"]/text()7
  \end{enumerate}


\begin{solution}
\begin{enumerate}
  \item La liste des n{\oe}uds \texttt{ingredient}, avec leur n{\oe}ud-texte fils.
  \item Les attributs \texttt{unit}.
  \item Le texte indiquant la difficulté de la recette: ici, "moyenne".
  \end{enumerate}
\end{solution}

\subsubsection*{XPath: expression}

On suppose que la base documentaire "Recettes" se représente sous la forme d'un très grand document 
XML avec un élément-racine \texttt{recipes}, contenant des sous-éléments \texttt{recipe} comme dans
l'exemple précédent. 
\begin{verbatim}
<?xml version="1.0" encoding="UTF-8"?>
<recipes>
  <recipe measures="FR">
    (... une recette...)
  </recipe>
  <recipe measures="US">
    (... une autre ...)
  </recipe>
</recipes>
\end{verbatim}

Donnez en XPath les expressions permettant de récupérer les informations suivantes:

\begin{enumerate}
\item Titre des recettes contenant du chocolat.
\item Quantité  de farine dans la recette des gougères.
\item Combien d'ingrédients dans la recette des profiteroles?
\item Titre des recettes pour lesquelles on n'a pas d'instructions.
\end{enumerate}


\begin{solution}
\begin{enumerate}
  \item \verb7//recipe[ingredient_lines/ingredient_line/ingredient="chocolat"]/title7
  \item \verb7//recipe[title="gougères"]/ingredient_lines/ingredient_line[ingredient/text()="farine"]/quantity7
   \item \verb7count(//recipe[title="profiteroles"]/ingredient_lines/ingredient_line)7
   \item \verb7//recipe[not(instructions)]/title7
  \end{enumerate}
\end{solution}

\section*{Seconde partie : indexation (5 pts)}

Nous souhaitons indexer les recettes. On va s'intéresser en particulier au champ \texttt{instructions}.
Voici quelques exemples:

La \emph{panna cotta} (identifiant '\texttt{pc}'):
\begin{small}
\begin{verbatim}
Mettre la crème, le sucre et la vanille dans une casserole et faire frémir. Ajouter 
les 3 feuilles de gélatine préalablement  trempées dans l'eau froide. Bien remuer 
et verser la crème dans des coupelles. Laisser refroidir quelques heures.
\end{verbatim}
\end{small}
la \emph{crème brulée} (identifiant 'cb')
\begin{small}
\begin{verbatim}
Faire bouillir le lait, ajouter la crème et le sucre hors du feu. Ajouter les jaunes 
d'oeufs, mettre au four au bain marie et laisser cuire doucement 
à 180C environ 10 minutes. Laisser refroidir puis mettre dessus du sucre roux et le 
brûler avec un petit chalumeau.
\end{verbatim}
\end{small}

et la \emph{mousse au chocolat} (identifiant 'mc').
\begin{small}
\begin{verbatim}
Faire ramollir le chocolat dans une terrine. Incorporer les jaunes et le sucre.
Puis, battre les blancs en neige ferme et les ajouter délicatement au 
mélange à l'aide d'une spatule. Mettre au frais 1 heure ou 2 minimum.
\end{verbatim}
\end{small}


\begin{enumerate}
  \item Cr\'eer les listes inverses pour les mots suivants : ''crème, sucre, oeuf, gélatine''.\\
  Vous indiquerez quel traitement préalable vous effectuez
  sur le texte, le nombre de termes que vous prenez en compte dans chaque document, l'idf des
  termes, et le tf dans chaque liste. 
\begin{solution}
IDF: 
\begin{itemize}
  \item \emph{crème} apparaît dans 2 documents sur 3, idf=3/2; 
   \item \emph{sucre} dans 3 sur 3 (idf=1), 
   \item \emph{oeuf} dans
1 sur 3 (idf=3) comme \emph{gélatine}.
\end{itemize}
Traitement du texte: on retire les articles, prépositions (stop words), les pluriels des noms. 
Le document pc contient 24 termes; le document cb contient 29 termes; le document mc 20 termes.
\begin{itemize}
  \item crème : pc (2/24), cb (1/29)
  \item sucre : cb (2/29), mc (1/20), pc (1/24)
  \item oeuf : cb (1/29)
  \item gélatine: pc (1/24)
\end{itemize}
\end{solution}

  \item Donner les résultats classés par combinaison tf.idf pour les requêtes suivantes (pas la peine d'applique le logarithme).
  
  \begin{itemize}
    \item crème et sucre
    \item crème et oeuf
    \item oeuf et gélatine
  \end{itemize}
\begin{solution}
\begin{itemize}
  \item crème et sucre: cb (3/2 * 1/29 + 1*2/29), pc (3/2*2/24 + 1*1/24), mc (1/20)
  \item crème et oeuf: cb (3/2 * 1/29 + 3 * 1/29), pc (3/2 * 2/24) 
  \item oeuf et gélatine: pc (3*1/24), cb (3*1/29)
\end{itemize}
\end{solution}  
  \item Commentez le résultat de la dernière requête. Est-il correct intuitivement? Que penser de l'indexation
    du terme 'oeuf', est-elle réprésentative du contenu des recettes?
  \end{enumerate}

\begin{solution}
Clairement, il faut des oeufs dans la crème au chocolat, mais le mot lui-même n'apparaît pas, ce qui fausse
les résultats. The \texttt{tf} est à 0 pour le document mc, et l'idf de 'oeuf' est surévalué (parce que le
nombre de documents lui-même est ridiculement petit).
\end{solution}

\section*{Troisième partie : systèmes NoSQL (8 pts)}

\begin{enumerate}
  \item \textbf{Réplication}: rappeler la définition des réplications synchrones et asynchrones, et indiquez
      brièvement les avantages/inconvénients de chacune.
      
  \item \textbf{Architecture}. Dans quelle architecture peut-on aboutir à des écritures conflictuelles? Donnez un exemple.
  
  \item \textbf{Gestion des pannes}. Rappeler la règle du quorum (majorité des votants) en cas de partitionnement de réseau, et
        justifiez-là.
        
   \item \textbf{MapReduce}. Dans une exécution MapReduce, le \textit{reducer} peut-il commencer à s'exécuter
      avant la fin de la phase de \textit{map}? Expliquer pourquoi (vous pouvez donner un exemple).
\end{enumerate}

\end{document}
